% GENERATED FILE. Edit canonical lessons, then run npm run generate:books.
% canonical-source-hash: fnv1a64:3c04f0d217cdb87d
% canonical-lessons: ES-C427-vendedor, ES-C427-comprador, ES-C427-policia, ES-C427-bombero, ES-C427-abogado, ES-C427-empresa, ES-C427-repaso-oficios, ES-C427-sintesis-profesion

\chapter{People Named By What They Do}
\label{ch:oficios}
\hlchaptermodality{427}

\begin{chapteropening}
\textbf{By the end of this chapter:} \emph{I can name five professions and the company somebody works for, take any profession word apart to see how it was built, and answer a question about my work without putting an article where Spanish leaves one out.}

Answer the question about what somebody does for a living, naming the profession with no article and the company with one.
\end{chapteropening}

\section[el vendedor]{el vendedor — the ending you already have, doing its real job}

\label{lesson:ES-C427-vendedor}



\begin{quote}
Say \emph{vender}. Say \emph{trabajador}. You have the verb and you have the ending. Put them together before reading on.
\end{quote}

\begin{cousinweb}
\textbf{el vendedor} — \textquotedblleft{}the seller, the shop assistant.\textquotedblright{}

Take \emph{vender}, drop the \textbf{-er}, add \textbf{-dor}. That is the whole operation, and you met the ending as an adjective in the character chapter. Here it does what it was built for: \textbf{making a person out of an action.}

\noindent
\begin{tabularx}{\linewidth}{@{}>{\raggedright\arraybackslash}X>{\raggedright\arraybackslash}X@{}}
\toprule
\textbf{verb} & \textbf{the one who does it} \\
\midrule
trabajar & el trabaja\textbf{dor} \\
vender & el vende\textbf{dor} \\
jugar & el juga\textbf{dor} \\
ganar & el gana\textbf{dor} \\
\bottomrule
\end{tabularx}

The ending is \textbf{productive}, which is the property worth having. You are not learning a word; you are learning a machine that turns any verb you know into the person who does it.
\end{cousinweb}

\begin{grammarlens}[title={Grammar Lens: where the ending puts the stress}]
Every \textbf{-dor} noun stresses its last syllable: \emph{vende-DOR}, \emph{trabaja-DOR}, \emph{juga-DOR}. None of them carries a written accent, because a word ending in \textbf{-r} is stressed on the last syllable by default and the accent marks only departures from the default.

The feminine adds \textbf{-a} and nothing else: \emph{la vendedora}, \emph{la trabajadora}.
\end{grammarlens}

\subsection*{Your turn}
Say these aloud:

\begin{itemize}
\raggedright
  \item \textquotedblleft{}el vendedor\textquotedblright{} — \emph{vende-DOR}
  \item the feminine form
  \item the two steps that build it from \emph{vender}
  \item the person-word for somebody who plays, built the same way
\end{itemize}

\begin{tcolorbox}[breakable,colback=teal!4,colframe=teal!35!black,title={Before you move on}]
What does the \textbf{-dor} ending do to a verb? (Makes a person out of it.) Why does \emph{vendedor} carry no written accent? (A word ending in \textbf{-r} is stressed on the last syllable anyway.)
\end{tcolorbox}
\section[el comprador]{el comprador — the other side of the counter}

\label{lesson:ES-C427-comprador}



\begin{quote}
Say \emph{el vendedor}. Say \emph{comprar}. Build the person who does the second thing, using the machine from the first.
\end{quote}

\begin{cousinweb}
\textbf{el comprador} — \textquotedblleft{}the buyer.\textquotedblright{}

\emph{Comprar} minus \textbf{-ar}, plus \textbf{-dor}. No new rule, which is the point of teaching the pair together: one ending buys both halves of a transaction.

\noindent
\begin{tabularx}{\linewidth}{@{}>{\raggedright\arraybackslash}X>{\raggedright\arraybackslash}X>{\raggedright\arraybackslash}X@{}}
\toprule
\textbf{side} & \textbf{verb} & \textbf{person} \\
\midrule
sells & vender & el vendedor \\
buys & comprar & el comprador \\
\bottomrule
\end{tabularx}
\end{cousinweb}

\begin{grammarlens}[title={Grammar Lens: the pair that names an economy}]
These two words are how a form, a contract or a news item names the parties to a sale. They are not shop vocabulary — \emph{el vendedor} in a shop is the assistant, but \emph{vendedor} and \emph{comprador} together are the \textbf{roles}, the way English says \emph{seller} and \emph{buyer} rather than \emph{shopkeeper} and \emph{customer}.

That distinction is why the source lists them under commerce rather than under shopping, and why knowing \emph{comprar} was not enough to cover it.
\end{grammarlens}

\subsection*{Your turn}
Say these aloud:

\begin{itemize}
\raggedright
  \item \textquotedblleft{}el comprador\textquotedblright{} — \emph{compra-DOR}
  \item the pair, seller then buyer
  \item both in the feminine
  \item which of the two you are when you hand over the money
\end{itemize}

\begin{tcolorbox}[breakable,colback=teal!4,colframe=teal!35!black,title={Before you move on}]
What two steps build \emph{comprador} from \emph{comprar}? (Drop the \textbf{-ar}, add \textbf{-dor}.) Are \emph{vendedor} and \emph{comprador} shop words or role words? (Roles — the two parties to a sale.)
\end{tcolorbox}
\section[el policía]{el policía — where the article changes the meaning}

\label{lesson:ES-C427-policia}



\begin{quote}
Say \emph{el comprador} and \emph{la compradora}, where the article follows the person. The word in this lesson does something different, and it is worth noticing.
\end{quote}

\begin{cousinweb}
\textbf{el policía} — \textquotedblleft{}the police officer.\textquotedblright{} Greek \textbf{pólis}, \textquotedblleft{}the city,\textquotedblright{} through Latin \emph{politia}, \textquotedblleft{}the administration of a state.\textquotedblright{}

The root is unusually productive, and every word on it is about the city:

\noindent
\begin{tabularx}{\linewidth}{@{}>{\raggedright\arraybackslash}X>{\raggedright\arraybackslash}X@{}}
\toprule
\textbf{word} & \textbf{the city in it} \\
\midrule
\textbf{polít}ica & the running of the city \\
metró\textbf{polis} & the mother-city \\
\textbf{polic}ía & the keeping of the city's order \\
\bottomrule
\end{tabularx}

So a police officer is named after \textbf{the city}, not after policing. The institution came first and the person was named from it.
\end{cousinweb}

\begin{grammarlens}[title={Grammar Lens: the article does the distinguishing}]
This is the fact a reading paper can test, and it catches people:

\noindent
\begin{tabularx}{\linewidth}{@{}>{\raggedright\arraybackslash}X>{\raggedright\arraybackslash}X@{}}
\toprule
\textbf{} & \textbf{meaning} \\
\midrule
\textbf{el} policía & a police officer, a man \\
\textbf{la} policía & the police \textbf{force} — or a woman officer \\
\bottomrule
\end{tabularx}

The word does not change. \textbf{The article carries the whole difference}, and \emph{la policía} is ambiguous between the force and a female officer in a way that only context settles.

Compare it with the \textbf{-dor} pair from the last two lessons, where the ending moved and the article merely followed. Here the ending is fixed at \textbf{-ía} and the article is doing all the work.
\end{grammarlens}

\subsection*{Your turn}
Say these aloud:

\begin{itemize}
\raggedright
  \item \textquotedblleft{}el policía\textquotedblright{} — \emph{po-li-CÍ-a}, four syllables
  \item what \emph{la policía} means, and why it has two readings
  \item which Greek word is inside it, and what it meant
  \item how this differs from \emph{el vendedor} / \emph{la vendedora}
\end{itemize}

\begin{tcolorbox}[breakable,colback=teal!4,colframe=teal!35!black,title={Before you move on}]
What does Greek \emph{pólis} mean? (The city.) What is the difference between \emph{el policía} and \emph{la policía}? (The first is an officer; the second is the force, or a woman officer.)
\end{tcolorbox}
\section[el bombero]{el bombero — named after the pump, not the fire}

\label{lesson:ES-C427-bombero}



\begin{quote}
Say \emph{el policía}. The other emergency service has a name built a different way again — and it is named after a piece of equipment.
\end{quote}

\begin{cousinweb}
\textbf{el bombero} — \textquotedblleft{}the firefighter.\textquotedblright{}

From \textbf{la bomba}, which in Spanish means \textbf{a pump} as well as a bomb. A \emph{bombero} is the person who works the pump.

English names the same person after what they fight: \emph{fire}-fighter, \emph{fire}-man. Spanish names them after what they carry. Two languages looking at the same person and choosing opposite halves of the scene.

The \textbf{-ero} ending is the third person-making ending in this chapter, and it attaches to \textbf{things} rather than to verbs:

\noindent
\begin{tabularx}{\linewidth}{@{}>{\raggedright\arraybackslash}X>{\raggedright\arraybackslash}X@{}}
\toprule
\textbf{thing} & \textbf{the person who works with it} \\
\midrule
la bomba — the pump & el bomb\textbf{ero} \\
el pan — bread & el pan\textbf{adero} \\
la carta — the letter & el cart\textbf{ero} \\
\bottomrule
\end{tabularx}

So Spanish has \textbf{-dor} for \emph{the one who does the action} and \textbf{-ero} for \emph{the one who works with the thing}.
\end{cousinweb}

\begin{grammarlens}[title={Grammar Lens: three endings, three logics}]
\noindent
\begin{tabularx}{\linewidth}{@{}>{\raggedright\arraybackslash}X>{\raggedright\arraybackslash}X>{\raggedright\arraybackslash}X@{}}
\toprule
\textbf{ending} & \textbf{built from} & \textbf{example} \\
\midrule
\textbf{-dor} & a verb & vendedor \\
\textbf{-ero} & a thing & bombero \\
\textbf{-ía} & an institution & policía \\
\bottomrule
\end{tabularx}

None of the three is guessable from the meaning; all three are readable once you know which is which. That is the honest position: you cannot predict which ending a profession takes, but you can take any profession apart.
\end{grammarlens}

\subsection*{Your turn}
Say these aloud:

\begin{itemize}
\raggedright
  \item \textquotedblleft{}el bombero\textquotedblright{} — \emph{bom-BE-ro}
  \item what \emph{la bomba} means besides a bomb
  \item the feminine form (\emph{la bombera})
  \item the difference between what \textbf{-dor} and \textbf{-ero} attach to
\end{itemize}

\begin{tcolorbox}[breakable,colback=teal!4,colframe=teal!35!black,title={Before you move on}]
What is a \emph{bombero} named after? (The pump.) What does \textbf{-ero} attach to, as against \textbf{-dor}? (A thing, where \textbf{-dor} attaches to a verb.)
\end{tcolorbox}
\section[el abogado]{el abogado — somebody called to your side}

\label{lesson:ES-C427-abogado}



\begin{quote}
Say \emph{el bombero}, named after his pump. The next profession is named after something nobody carries: a summons.
\end{quote}

\begin{cousinweb}
\textbf{el abogado} — \textquotedblleft{}the lawyer.\textquotedblright{}

Latin \textbf{advocatus}: \textbf{ad-} to + \textbf{vocatus}, \textquotedblleft{}called.\textquotedblright{} Literally \textbf{one called to your side} — in a Roman court, the friend you summoned to speak for you.

English borrowed the same word twice and kept both:

\noindent
\begin{tabularx}{\linewidth}{@{}>{\raggedright\arraybackslash}X>{\raggedright\arraybackslash}X@{}}
\toprule
\textbf{word} & \textbf{the calling in it} \\
\midrule
\textbf{advocate} & one called to your side \\
\textbf{vocation} & a calling, in the older sense \\
\textbf{convoke} & to call together \\
\bottomrule
\end{tabularx}

\emph{Vocare}, \textquotedblleft{}to call,\textquotedblright{} is the root under all of them, and it is the same root as \emph{voz}, the voice.

Spanish wore \emph{advocatus} down to \emph{abogado} while English kept the Latin shape intact, which is the ordinary difference between a word that was \textbf{inherited} and a word that was \textbf{borrowed} — the same distinction \emph{actor} and \emph{trabajador} demonstrated.
\end{cousinweb}

\begin{grammarlens}[title={Grammar Lens: the profession without a visible ending}]
\noindent
\begin{tabularx}{\linewidth}{@{}>{\raggedright\arraybackslash}X>{\raggedright\arraybackslash}X@{}}
\toprule
\textbf{profession} & \textbf{how the word was made} \\
\midrule
vendedor & Spanish added \textbf{-dor} to a Spanish verb \\
bombero & Spanish added \textbf{-ero} to a Spanish noun \\
abogado & Latin made it, and Spanish inherited it whole \\
\bottomrule
\end{tabularx}

\emph{Abogado} has no ending you can remove to find a Spanish verb, because it was never built in Spanish. It arrived finished.

The feminine is the ordinary swap: \emph{la abogada}.
\end{grammarlens}

\subsection*{Your turn}
Say these aloud:

\begin{itemize}
\raggedright
  \item \textquotedblleft{}el abogado\textquotedblright{} — \emph{a-bo-GA-do}
  \item the feminine form
  \item what the two Latin pieces mean (to, and called)
  \item which English word is the same Latin word kept whole (\emph{advocate})
\end{itemize}

\begin{tcolorbox}[breakable,colback=teal!4,colframe=teal!35!black,title={Before you move on}]
What did \emph{advocatus} literally mean? (One called to your side.) Can you find a Spanish verb inside \emph{abogado}? (No — it was inherited whole from Latin rather than built in Spanish.)
\end{tcolorbox}
\section[la empresa]{la empresa — a thing taken on}

\label{lesson:ES-C427-empresa}



\begin{quote}
Say \emph{el vendedor} and \emph{el comprador}. You can name both parties to a sale. Say where the seller works, and notice you have no word for it.
\end{quote}

\begin{cousinweb}
\textbf{la empresa} — \textquotedblleft{}the company, the firm.\textquotedblright{}

From Latin \textbf{in-} + \textbf{prehendere}, \textquotedblleft{}to take hold of.\textquotedblright{} An \emph{empresa} is literally \textbf{a thing taken on} — an undertaking.

English took the same idea through French and got \textbf{enterprise}, from \emph{entre-} + \emph{prendre}, the same two pieces in the other Romance language. The words are cousins built to the same plan.

The difference is what each language did with it afterwards. English \emph{enterprise} mostly means the \textbf{venture}; Spanish \emph{empresa} means both the venture \textbf{and the firm that carries it out}. For an A1 form asking where you work, \emph{empresa} is the everyday word — it is what English calls \emph{the company}.

\emph{Prehendere} is generous, and two of its descendants are already useful to you: \emph{aprender}, to learn, is to \textbf{take hold of} something; and \emph{sorprender}, to surprise, is to take somebody from above.
\end{cousinweb}

\begin{grammarlens}[title={Grammar Lens: a workplace rather than a worker}]
The chapter has named five people by what they do. \textbf{La empresa} is the one word in it that names a \textbf{place and an organisation} instead:

\noindent
\begin{tabularx}{\linewidth}{@{}>{\raggedright\arraybackslash}X>{\raggedright\arraybackslash}X@{}}
\toprule
\textbf{word} & \textbf{what it names} \\
\midrule
el vendedor & a person \\
el abogado & a person \\
la empresa & the organisation they work for \\
\bottomrule
\end{tabularx}

That is why a form has two separate lines, \emph{profesión} and \emph{empresa}, and why answering one does not answer the other.
\end{grammarlens}

\subsection*{Your turn}
Say these aloud:

\begin{itemize}
\raggedright
  \item \textquotedblleft{}la empresa\textquotedblright{} — \emph{em-PRE-sa}
  \item what the two Latin pieces mean (in, and take hold of)
  \item the English cousin, and the one difference in what it covers
  \item which verb for learning comes from the same root (\emph{aprender})
\end{itemize}

\begin{tcolorbox}[breakable,colback=teal!4,colframe=teal!35!black,title={Before you move on}]
What does \emph{empresa} literally describe? (A thing taken on.) Does it mean the venture or the firm? (Both — where English \emph{enterprise} mostly means only the venture.)
\end{tcolorbox}
\section[cinco personas y un lugar]{cinco personas y un lugar — sorted by how each name was built}

\label{lesson:ES-C427-repaso-oficios}



\begin{quote}
Name as many of the six as you can before reading the table.
\end{quote}

\begin{grammarlens}[title={Grammar Lens: four ways to name a person}]
\noindent
\begin{tabularx}{\linewidth}{@{}>{\raggedright\arraybackslash}X>{\raggedright\arraybackslash}X@{}}
\toprule
\textbf{word} & \textbf{built from} \\
\midrule
el vendedor & the verb \emph{vender}, plus \textbf{-dor} \\
el comprador & the verb \emph{comprar}, plus \textbf{-dor} \\
el bombero & the noun \emph{bomba}, plus \textbf{-ero} \\
el policía & the institution, \emph{la policía} \\
el abogado & nothing in Spanish — inherited whole from Latin \\
\bottomrule
\end{tabularx}

Four different methods for five words, which is the honest summary: Spanish does not have one way of naming a profession. What it has is a small set of methods, and knowing them lets you \textbf{take a profession apart} even when you could not have predicted it.

The \textbf{-dor} pair is the only one that is productive for you: give it any verb you know and it gives you back the person — which is how \emph{trabajar} gave \emph{trabajador} in the character chapter, before you had any use for the ending.
\end{grammarlens}

\begin{grammarlens}[title={Grammar Lens: which ones change, and how}]
\noindent
\begin{tabularx}{\linewidth}{@{}>{\raggedright\arraybackslash}X>{\raggedright\arraybackslash}X@{}}
\toprule
\textbf{masculine} & \textbf{feminine} \\
\midrule
el vendedor & la vendedora \\
el bombero & la bombera \\
el abogado & la abogada \\
\textbf{el policía} & \textbf{la policía} — and the word also means the force \\
\bottomrule
\end{tabularx}

Three behave ordinarily. \textbf{Policía} is the exception worth remembering, because the article is carrying a distinction no ending marks.

And \textbf{la empresa} is not a person at all — it is where several of these people work, which is why a form asks for it on its own line.
\end{grammarlens}

\subsection*{Your turn}
Say these aloud:

\begin{itemize}
\raggedright
  \item the two built with \textbf{-dor}, and the verb inside each
  \item the one built with \textbf{-ero}, and the thing inside it
  \item the one whose article changes the meaning, and both readings
  \item the one that is not a person
\end{itemize}

\begin{tcolorbox}[breakable,colback=teal!4,colframe=teal!35!black,title={Before you move on}]
Which of the endings can you apply to a new verb yourself? (\textbf{-dor}.) Which word's meaning is decided by its article? (\emph{Policía}.) Which of the six is not a person? (\emph{La empresa}.)
\end{tcolorbox}
\section[¿A qué se dedica?]{¿A qué se dedica? — the question about work}

\label{lesson:ES-C427-sintesis-profesion}



\begin{quote}
Say \emph{el abogado} and \emph{la empresa}. One answers what somebody does and one answers where. A conversation about work needs both.
\end{quote}

\begin{grammarlens}[title={Grammar Lens: naming a profession takes no article}]
This is the one grammatical trap in the answer, and English speakers fall into it every time:

\noindent
\begin{tabularx}{\linewidth}{@{}>{\raggedright\arraybackslash}X>{\raggedright\arraybackslash}X@{}}
\toprule
\textbf{} & \textbf{} \\
\midrule
correct & Soy abogado. \\
wrong & Soy \textbf{un} abogado. \\
\bottomrule
\end{tabularx}

English says \emph{I am \textbf{a} lawyer}. Spanish drops the article when the sentence states nothing but the profession, because the bare noun is working as a description rather than picking out one individual.

The article comes back as soon as you add anything to it: \emph{Soy \textbf{un} abogado trabajador} — \textquotedblleft{}I am a hard-working lawyer.\textquotedblright{} Adding the adjective turns the description into a particular kind of thing, and a particular thing takes an article.
\end{grammarlens}

\begin{tcolorbox}[breakable,skin=enhanced,colback=orange!6,colframe=orange!45!black,boxrule=0.5pt,arc=1mm,left=6pt,right=6pt,top=4pt,bottom=4pt,fonttitle=\bfseries,title={Read this}]
Read aloud. Every word has been taught:

\begin{quote}
— ¿A qué se dedica? — Soy abogado. Trabajo en una empresa grande. — ¿Y su hermana? — Es policía. Mi hermano es bombero.
\end{quote}

Three professions, no article on any of them, and \emph{una empresa} taking one because it is a thing rather than a description.
\end{tcolorbox}

\subsection*{Your turn}
Say these aloud:

\begin{itemize}
\raggedright
  \item \textquotedblleft{}¿A qué se dedica?\textquotedblright{}
  \item \textquotedblleft{}Soy abogado\textquotedblright{}, and stop yourself before putting \emph{un} in
  \item the same sentence with an adjective added, where the article returns
  \item your own answer, with the profession and the company
\end{itemize}

\begin{tcolorbox}[breakable,colback=teal!4,colframe=teal!35!black,title={Before you move on}]
Why is it \emph{Soy abogado} and not \emph{Soy un abogado}? (A bare profession is a description, and a description takes no article.) When does the article come back? (As soon as an adjective is added.)
\end{tcolorbox}
